%-----------------------------------------------------------------
% Rodrigo Chang, 2019
% Banco de Guatemala
% Programa de Estudios Superiores 2019-2020
% Diapositivas del curso "Programación I"
%
% Descargo de responsabilidad.
%   El contenido de este documento es responsabilidad exclusiva de
% su(s) autores. Cualquier opinión expresada no necesariamente
% representa la postura institucional del Banco de Guatemala. 
% 
% Contenido.
%   El contenido de estas diapositivas está construido a partir de 
% modificaciones al material provisto por Ana Bell, Eric Grimson y 
% John Guttag para el curso 6.0001 Introduction to Computer Science 
% and Programming in Python, otoño 2016, Massachusetts Institute of 
% Technology: MIT OpenCourseWare, https://ocw.mit.edu
%
% Licencia: 
% Creative Commons BY-NC-SA
% https://creativecommons.org/licenses/by-nc-sa/4.0/.
%
% Formato. 
%   El formato utilizado para la elaboración de este documento 
% ha sido modificado a partir del repositorio: 
%       https://github.com/rinnert/statlearning-python-pwf
%
% Kurt Rinnert, 2019
%   Copying and distribution of this file, with or without 
% modification, are permitted in any medium without royalty provided 
% the copyright notice and this notice are preserved. This file is 
% offered as-is, without any warranty.
%
%-----------------------------------------------------------------
\documentclass[mathserif]{beamer}
% Para configurar las diapositivas en relación 16:9, 
% agregar la opción aspectratio=169

% Título de la presentación y tópicos a tratar
\usepackage{arev, arevmath}
\usepackage[scaled]{cabin}
%\usepackage[T1]{fontenc}
\usepackage[super]{nth}
\usepackage{pifont}
\usepackage{wasysym}
\usepackage{tabularx}
\usepackage{array}
\usepackage{booktabs}
\usepackage{boldline}
\usepackage{colortbl}
\usepackage{amsmath}
\usepackage{bm}
\usepackage{tcolorbox}
\usepackage{adjustbox}
\usepackage{minibox}
\usepackage{makecell}
\usepackage{adjustbox}
\usepackage{textcomp}
\usepackage[absolute,overlay]{textpos}
\setlength{\TPHorizModule}{1mm}%
\setlength{\TPVertModule}{1mm}%
\tcbuselibrary{skins}
\usepackage{graphicx}
\usepackage[spanish, es-tabla]{babel}
\usepackage[utf8]{inputenc}
\usepackage{csquotes}

\usepackage{comment}
\usepackage{svg}

\usepackage{verbatim}
\usepackage{pdfpages}

% Columnas múltiples
\usepackage{multicol}

% -------------------------------------------
%	LISTINGS
% -------------------------------------------
% Paquetes necesarios para los códigos
\usepackage{listings}
\usepackage{color}
% Configuración de listings (pendiente)
%\input{conf/listingStyles}
% Redefinición de titulos de caption
\renewcommand\lstlistingname{\textbf{Código}}
\renewcommand\lstlistlistingname{Códigos}

% Definición de colores para Python (Jupyter notebooks)
\definecolor{python_keyword_color}{RGB}{0, 128, 0}
\definecolor{python_string_color}{RGB}{186, 34, 34}
\definecolor{python_comment_color}{RGB}{1, 121, 121}
\definecolor{frame_numbers_color}{rgb}{0.5,0.5,0.5}
\definecolor{python_identifier_color}{RGB}{0,0,0}
\definecolor{jupyter_background}{RGB}{247, 247, 247}
\definecolor{jupyter_frame}{RGB}{207, 207, 207}

% Definicion de estilo para Python
\lstdefinestyle{custom_python}{
	language=python,			% the language of the code
	backgroundcolor=\color{jupyter_background}, 
	% the size of the fonts that are used for the code
	basicstyle=\linespread{1}\small\ttfamily,	
	% Configuración de colores
	commentstyle=\color{python_comment_color},	% comment style
	identifierstyle=\color{python_identifier_color},
	keywordstyle=\mdseries\color{python_keyword_color},	% keyword style
	numberstyle=\tiny\color{frame_numbers_color},	
	stringstyle=\color{python_string_color}, % string literal style 
	captionpos=t,				% sets the caption-position to bottom
	breakatwhitespace=false,	% sets if automatic breaks should only happen at whitespace
	breaklines=true,			% sets automatic line breaking
	morekeywords={},			% if you want to add more keywords to the set
	deletekeywords={eps},		% delete keywords from the given language
	escapeinside={\%*}{*)},		% if you want to add LaTeX within your code
	framesep=3pt,
	keepspaces=true,			% keeps spaces in text,       
	numbers=none, 	% (none, left, right)
	numbersep=5pt,	% how far the line-numbers are from the code
	% if not set, the frame-color may be changed on line-breaks within not-black text (e.g. comments (green here))
	rulecolor=\color{jupyter_frame}, 
	% show spaces everywhere adding particular underscores; it overrides 'showstringspaces'
	showspaces=false,                
	% underline spaces within strings only
	showstringspaces=false,          
	% show tabs within strings adding particular underscores
	showtabs=false,                  
	stepnumber=0,	% the step between two line-numbers      
	% sets default tabsize to 2 spaces
	tabsize=2,                       
	% lets you use non-ASCII characters; for 8-bits encodings only, does not work with UTF-8
	extendedchars=true,      
	% Support additional characters            
	literate=
  {á}{{\'a}}1 {é}{{\'e}}1 {í}{{\'i}}1 {ó}{{\'o}}1 {ú}{{\'u}}1
  {Á}{{\'A}}1 {É}{{\'E}}1 {Í}{{\'I}}1 {Ó}{{\'O}}1 {Ú}{{\'U}}1
  {à}{{\`a}}1 {è}{{\`e}}1 {ì}{{\`i}}1 {ò}{{\`o}}1 {ù}{{\`u}}1
  {À}{{\`A}}1 {È}{{\'E}}1 {Ì}{{\`I}}1 {Ò}{{\`O}}1 {Ù}{{\`U}}1
  {ä}{{\"a}}1 {ë}{{\"e}}1 {ï}{{\"i}}1 {ö}{{\"o}}1 {ü}{{\"u}}1
  {Ä}{{\"A}}1 {Ë}{{\"E}}1 {Ï}{{\"I}}1 {Ö}{{\"O}}1 {Ü}{{\"U}}1
  {â}{{\^a}}1 {ê}{{\^e}}1 {î}{{\^i}}1 {ô}{{\^o}}1 {û}{{\^u}}1
  {Â}{{\^A}}1 {Ê}{{\^E}}1 {Î}{{\^I}}1 {Ô}{{\^O}}1 {Û}{{\^U}}1
  {œ}{{\oe}}1 {Œ}{{\OE}}1 {æ}{{\ae}}1 {Æ}{{\AE}}1 {ß}{{\ss}}1
  {ű}{{\H{u}}}1 {Ű}{{\H{U}}}1 {ő}{{\H{o}}}1 {Ő}{{\H{O}}}1
  {ç}{{\c c}}1 {Ç}{{\c C}}1 {ø}{{\o}}1 {å}{{\r a}}1 {Å}{{\r A}}1
  {€}{{\EUR}}1 {£}{{\pounds}}1,   
}

% Configurar el listing para el documento actual
\lstset{
	style = custom_python, 
	frame = single
}

% -------------------------------------------
%	OTRAS CONFIGURACIONES
% -------------------------------------------
\makeatletter
\newcommand{\antsize}{\@setfontsize{\antsize}{4pt}{4pt}}
\makeatother
\newcommand{\at}{\makeatletter @\makeatother}

\newcommand{\cmark}{\ding{51}}%
\newcommand{\bottomline}[1]{\vskip0pt plus 1fill{\alert{#1}\phantom{g}\vskip 1.0mm}}

% Rutas para las imágenes
%\graphicspath{{figs/}{figs/common/}}

% Separación de filas en las tablas más amplia
\renewcommand{\arraystretch}{1.2}


% -------------------------------------------
%	ESTILO DE LA PRESENTACIÓN
% -------------------------------------------

\mode<presentation>%
{
	\usetheme{default}
	%\usetheme[width=2.5cm]{PaloAlto}
	\usecolortheme{dove}
	\useoutertheme{infolines}
	% oder auch nicht

	% BANGUAT Colors
	\definecolor{BANGUATblue}{RGB}{37,86,162} % 0x255682
	\definecolor{BANGUATorange}{RGB}{255,130,0} % 0xff8200
	\definecolor{BANGUATgreen}{RGB}{0,100,0}
	\definecolor{BANGUATdark}{RGB}{80,80,80} % 0x505050
	\definecolor{BANGUATlight}{RGB}{120,120,120}
	\definecolor{BANGUATbrown}{RGB}{178,91,0}

	\definecolor{codebg}{rgb}{0.95,0.95,0.95}

	% Color theme
	\setbeamercolor{alerted text}{fg=BANGUATorange}
	\setbeamercolor{frametitle}{fg=BANGUATblue}
	\setbeamercolor{title}{fg=BANGUATblue}
	\setbeamercolor{subtitle}{fg=BANGUATorange}
	\setbeamercolor{normal text}{fg=BANGUATdark}
	\setbeamercolor{author in foot}{fg=BANGUATblue}
	\setbeamercolor{item}{fg=BANGUATblue}
	\setbeamercolor{footline}{fg=BANGUATblue}
	%\setbeamercolor{item projected}{bg=BANGUATorange}
	%\setbeamercolor{item projected}{fg=white}

	\setbeamertemplate{headline}
	{}
	\setbeamertemplate{frametitle}
	{
		%\textbf{{\insertframetitle\phantom{g}}}\\
		%\textbf{\insertframetitle\phantom{g}}\\
		\textbf{\underline{\insertframetitle\phantom{g}}}\\
		%\textbf{\underline{\insertframetitle}}\\
		\vskip 1.0mm
		%{\color{UOLgold}\hrule height 2pt}
		%\par
	}
	\addtobeamertemplate{frametitle}{}{\vspace{-1em}}
	\setbeamertemplate{footline}{
		{%
			\textbf{ \hskip 3.0mm\insertshorttitle\phantom{.}---\phantom{.}\insertshortinstitute\hfill\insertframenumber\,/\,\inserttotalframenumber\hskip 3.0mm} 
		}
	}

	\setbeamertemplate{navigation symbols}{}%remove navigation symbols
	\setbeamertemplate{itemize items}[circle]
	\setbeamertemplate{enumerate items}[fg=BANGUATblue]
	\setbeamercolor{itemize items}{fg=BANGUATblue}
	\setbeamercolor{sidebar}{bg=BANGUATblue}
	\setbeamercolor{title in sidebar}{fg=BANGUATorange}
	\setbeamercolor{author in sidebar}{fg=BANGUATorange}
	\setbeamercolor{section in sidebar}{fg=BANGUATorange}
}



% -------------------------------------------
%	CONFIGURACIÓN DE DIAPOSITIVA DE TÍTULO
% -------------------------------------------

% Título del curso
\title[Programación I]{\textbf{%
Introducción a las ciencias de la computación\\%
{\it y programación en Python}}\\%
%\href{https://ocw.mit.edu/courses/electrical-engineering-and-computer-science/6-0001-introduction-to-computer-science-and-programming-in-python-fall-2016/}{\tiny\it Basado en el curso ``Introduction to Computer Science and Programming in Python'' por Ana Bell, Eric Grimson, y John Guttag.}%
\vspace{3em}%
}
%\vspace{-2.5cm}{}

% Autor y correo electrónico
\author{\href{mailto:?to=Rodrigo Chang <rrcp@banguat.gob.gt>&subject=PES Programación I}{Rodrigo Chang}}

% Información del instituto
\institute[\href{https://www.banguat.gob.gt/}{Banco de Guatemala}]%
{\color{BANGUATblue} \bfseries \href{https://www.banguat.gob.gt/}{Banco de Guatemala}\\%
\vspace{3mm}%
\includegraphics[width=0.13\textwidth]{figs/common/logo-banguat.png}}

% Sin fecha en particular
\date{}

% Información de licencias y copyright
\titlegraphic{
	 %\texorpdfstring{\vspace{-3cm}}{}
	 \vspace{-1.5cm}
	 \begin{minipage}[b][1cm][b]{0.25\textwidth}%
	    \color{BANGUATlight}\antsize
		%Copyright \textcopyright~2019\\
		\href{mailto:?to=Rodrigo Chang <rrcp@banguat.gob.gt>&subject=PES Programación I}{Rodrigo Chang <rrcp@banguat.gob.gt>}\\
		Este material está construido a partir de modificaciones al material provisto por Ana Bell, Eric Grimson y John Guttag para el curso \textit{6.0001 Introduction to Computer Science and Programming in Python}, otoño 2016, Massachusetts Institute of Technology: MIT OpenCourseWare, \href{https://ocw.mit.edu}{https://ocw.mit.edu}. Licencia: \href{https://creativecommons.org/licenses/by-nc-sa/4.0/}{Creative Commons BY-NC-SA}.
		
		%Copying and distribution of this file, with or without modification, are permitted in any medium without royalty provided the copyright notice and this notice are preserved.  This file is offered as-is, without any warranty.


		%Some of the figures in this presentation are taken from ``An Introduction to Statistical Learning, with applications in R''  (Springer, 2013) with permission from the authors: G. James, D. Witten,  T. Hastie and R. Tibshirani 
	 \end{minipage}%
	 \hspace{9cm}
}

% Restar uno al contador
\addtocounter{framenumber}{-1}

% Cajas de colores
\newcommand{\tabboxset}{\tcbset{enhanced, nobeforeafter, boxrule=0pt, boxsep=0pt, colback=codebg, colframe=codebg, coltext=BANGUATdark, rounded corners, arc=4pt, fonttitle={\bfseries\tiny}}}
\newcommand{\codeboxset}{\tcbset{enhanced, nobeforeafter, boxrule=0pt, boxsep=0pt, colback=codebg, colframe=codebg, coltext=BANGUATdark, rounded corners, arc=4pt, fonttitle={\bfseries\tiny}}}

% Comandos útiles
\newcommand{\orange}{\color{BANGUATorange}}
\newcommand{\blue}{\color{BANGUATblue}}
\newcommand{\dark}{\color{BANGUATdark}}
\newcommand{\R}{\mathbb{R}}
\newcommand{\dat}[1]{{\footnotesize\tt\orange #1}}
\newcommand{\mdat}[1]{\text{\dat{#1}}}
\newcommand{\val}[1]{{\footnotesize\tt\blue #1}}
\newcommand{\mval}[1]{\text{\val{#1}}}
\newcommand{\e}[1]{\emph{#1}}
\newcommand{\bh}{\hat{\beta}}
\newcommand{\h}{\hat}
\newcommand{\bl}{\bottomline}


\makeatletter
\newcommand{\includegraphicsdpi}[3]{%
	\pdfimageresolution=#1%
	\includegraphics[#2]{#3}%
	\pdfimageresolution=72%
}

\newcommand{\Quote}[2]{%
	\begin{center} 
		\begin{minipage}{0.7\textwidth} 
			\hrule
			\vskip 3mm
			\emph{{\color{BANGUATblue} ``#1''}\\
			{\color{BANGUATorange} --- #2}}
			\vskip 3mm
			\hrule
			\vskip 2mm
		\end{minipage}
	\end{center}
}

% Texto centrado en la diapositiva, quizás para el abstract o para indicar abajo de un quote
\newenvironment{blurb}%
	{\begin{center}\begin{minipage}{0.6\textwidth}\footnotesize}
	{\end{minipage}\end{center}}

% Environment para códigos o algoritmos. Parecido a un bottomline, se puede utilizar para dar definiciones o enfatizar en algún concepto
\newenvironment{cpage}%
	{\begin{center}\begin{minipage}{0.75\textwidth}}
	{\end{minipage}\end{center}}

% Bloque (similar a un pop-up) de alerta para indicar resultados en tablas, definiciones, fórmulas, etc.
% Primer argumento es el tamaño 
% Segundo argumento es el título del bloque
\newenvironment{popblock}[2]%
	{\begin{center}\begin{minipage}{#1}\footnotesize
		\begin{tcolorbox}[colframe=codebg, colback=white, colupper=BANGUATdark, title={\normalsize\bfseries\blue #2}]}
	{\end{tcolorbox}\end{minipage}\end{center}}
\makeatother

% Incluir diapositivas en PDF
% 1 - Título de la diapositiva
% 2 - páginas del PDF a incluir
% 3 - ruta hacia el archivo PDF
\newcommand{\includeslides}[3]{%
{%
\setbeamercolor{background canvas}{bg=}%
\addtocounter{framenumber}{1}%
\includepdf[pages={#1}, link=true, pagecommand={%
\Large\blue\bfseries\underline{%
#3\phantom{g}}%
\vspace{-1em}%
\vskip 1.0mm}]{#2}%
}}
\subtitle{%
  \bfseries{Pruebas, depuración y manejo de excepciones}\\%
  %{\tiny\it Manejo de strings, \textit{guess and check}, métodos de aproximación}\\%
}

% Inicio del documento
\begin{document}
\frame[plain]{
	%\vskip 1.0mm
	\titlepage
	%\vskip 1.0mm
}
%-----------------------------------------------------------------
% Empieza el contenido de la presentación
%-----------------------------------------------------------------

\begin{frame}{Abstract}
	\Quote{Always code as if the guy who ends up maintaining your code will be a violent psychopath who knows where you live.}{John Woods}
	
	% Commenting your code is like cleaning your bathroom — you never want to do it, but it really does create a more pleasant experience for you and your guests. - Ryan Campbell
    
	\begin{blurb}
		Veremos algunos conceptos y buenas prácticas para realizar pruebas y depuración.
		
		Veremos algunas sentencias de {\blue excepción} y manejo de errores.
		
	\end{blurb}
\end{frame}

\section{Depuración y pruebas}

\begin{frame}[fragile]{Alta calidad: una analogía}
\vspace{0.5cm}
Supongamos que queremos hacer una sopa, pero algunos {\blue\bfseries insectos} caen del techo. ¿Qué harías?
    \begin{columns}
        \begin{column}{0.4\textwidth}
            \begin{itemize}
                \item Revisar la sopa a ver si hay bichos.
                \begin{itemize}
                    \item Pruebas.
                \end{itemize}
                \item Mantener la tapa cerrada.
                \begin{itemize}
                    \item Programación defensiva.
                \end{itemize}
                \item Limpiar la cocina.
                \begin{itemize}
                    \item Eliminamos la fuente de los bichos.
                \end{itemize}
                \end{itemize}
            \end{column}
        \begin{column}{0.6\textwidth}
        \begin{figure}
            \centering
            \includegraphics[width=0.98\textwidth]{figs/depuracion/debugging.jpg}
        \end{figure}
        \end{column}
    \end{columns}
\end{frame}

\begin{frame}[fragile]{}
    \begin{itemize}
        \item {\blue\bfseries Programación defensiva}
        \begin{itemize}
            \item Escribir {\orange especificaciones} para las funciones.
            \item {\orange Modularizar} los programas.
            \item Revisar {\orange condiciones} sobre entradas y salidas (\textit{assertions}).
        \end{itemize}
        
        \item {\blue\bfseries Validación y pruebas}
        \begin{itemize}
            \item {\orange Comparar} pares de entradas y salidas para verificar las especificaciones.
            \item ¡No está funcionando!
            \item Pensar: ¿cómo puedo hacer que falle mi programa?
        \end{itemize}
        
        \item {\blue\bfseries Depuración (\textit{debugging})}
        \begin{itemize}
            \item {\orange Estudiar eventos} que han llevado a errores.
            \item Pensar: ¿por qué no está funcionando?
            \item ¿Cómo puedo arreglar mi programa?
        \end{itemize}
    \end{itemize}
\end{frame}

% Set up for easy testing
\begin{frame}[fragile]{Programación defensiva}
    \begin{itemize}
        \item Desde el {\orange\bfseries principio}, {\blue\bfseries diseñar} el código para facilitar esta parte.
        
        \item Diseñar el programa en {\blue módulos} que puedan ser probados y depurados individualmente.
        
        \item {\blue\bfseries Documentar las restricciones} sobre módulos:
        \begin{itemize}
            \item ¿Cuál se espera que sea la entrada?
            \item ¿Cuál se espera que sea la salida?
        \end{itemize}
        
        \item {\blue\bfseries Documentar supuestos} detrás del código.
    \end{itemize}
\end{frame}

\begin{frame}[fragile]{¿Cuándo se debe probar un código?}
    \begin{itemize}
        \item Asegurarse que el código {\blue\bfseries pueda ejecutarse}.
        \begin{itemize}
            \item Eliminar errores de sintaxis, semántica estática.
            \item El intérprete de Python y los editores ayudan en esta tarea.
        \end{itemize}
        
        \item Contar con un {\blue\bfseries conjunto de prueba}.
        \begin{itemize}
            \item Conjunto de entrada del cual conocemos la salida esperada.
        \end{itemize}
    \end{itemize}
\end{frame}

% Tipos de pruebas
\begin{frame}[fragile]{Clases de pruebas}
    \begin{itemize}
        \item {\blue\bfseries Pruebas de unidad}
        \begin{itemize}
            \item Validar cada pieza de un programa.
            \item {\orange Probar cada función} de forma separada.
        \end{itemize}
        
        \item {\blue\bfseries Pruebas de regresión}
        \begin{itemize}
            \item Añadir casos de pruebas conforme se encuentran \textit{bugs}.
            \item {\orange Validar} errores reintroducidos y previamente arreglados.
        \end{itemize}
        
        \item {\blue\bfseries Pruebas de integración}
        \begin{itemize}
            \item ¿Cómo unir todo el programa?, ¿funciona en conjunto?
            \item No se debe apresurar en este paso.
        \end{itemize}
    \end{itemize}
\end{frame}


% Tipos de pruebas
\begin{frame}[fragile]{Tipos de pruebas}
    \begin{itemize}
        \item {\blue\bfseries Intuición} acerca de los límites del programa.
\begin{lstlisting}
def is_bigger(x, y):
    ''' Assumes x and y are ints
    Returns True if y is less than x, else False '''
\end{lstlisting}
        \begin{itemize}
            \item ¿alguna partición natural para probar esta función?
        \end{itemize}
        
        \item Si no hay particiones naturales $\Rightarrow$ {\blue pruebas aleatorias}.
        \begin{itemize}
            \item Más pruebas $\Rightarrow$ mayor prob. de éxito. 
            \item Método no ideal.
        \end{itemize}
        
        \item {\orange\bfseries Prueba de caja negra}
        \begin{itemize}
            \item Explorar casos a través de la especificación.
        \end{itemize}
        
        \item {\orange\bfseries Prueba de caja de cristal}
        \begin{itemize}
            \item Explorar casos a través del conocimiento del código.
        \end{itemize}
    \end{itemize}
\end{frame}

% Pruebas de caja negra
\begin{frame}[fragile]{Pruebas de caja negra}
\begin{lstlisting}
def sqrt(x, eps):
    ''' Assumes x, eps floats, x >= 0, eps > 0
    Returns res such that x-eps <= res*res <= x+eps '''
\end{lstlisting}
    \begin{itemize}
        \item Diseñada {\orange sin conocimiento} del código.
        \item Puede ser realizada por alguien distinto al programador.
        \item Puede ser {\orange reutilizada} aún si la implementación cambia.
        \item Algunos caminos:
        \begin{itemize}
            \item Casos de prueba en particiones naturales.
            \item Considera {\orange condiciones límite}: listas vacías, listas \textit{singleton}, números muy grandes o muy pequeños, etc.
        \end{itemize}
    \end{itemize}
\end{frame}

% Black box testing example
\includeslides{9}{slides-60001/MIT6_0001F16_Lec7_mod.pdf}{Pruebas de caja negra}

% Pruebas de caja de cristal
\begin{frame}[fragile]{Pruebas de caja de cristal}
    \begin{itemize}
        \item Utilizamos {\orange el código} para construir los casos de prueba.
        \item Pueden ser {\orange completas} si se explora cada ramificación del código al menos una vez.
        \item Desventajas:
        \begin{itemize}
            \item Muchas iteraciones.
            \item Nos podemos saltar alguna rama.
        \end{itemize}
        
        \item Guía para diseñarlas:
        \begin{itemize}
            \item Revisar condicionales: ¿están todos los posibles casos?
            \item Revisar ciclos: ¿entra al ciclo?, ¿por qué se ejecuta una vez, o más de una vez?
            \item Tipos de datos: ¿Devuelve el tipo esperado?, ¿se tuvo cuidado al operar las variables?
        \end{itemize}
    \end{itemize}
\end{frame}

% Ejemplo de glass box testing
\begin{frame}[fragile]{}
\begin{lstlisting}
def abs(x):
    """ Assumes x is an int
    Returns x if x>=0 and -x otherwise """
    if x < -1:
        return -x
    else:
        return x
\end{lstlisting}
    \begin{itemize}
        \item Aún si la prueba es {\orange completa}, podríamos saltarnos un \textit{bug}.
        \item ¿Qué pasa con los casos 2 y -2?
        \item ¿Qué pasa con \lstinline|abs(-1)|?
        \item Se deben probar casos {\orange\bfseries límite}.
    \end{itemize}
\end{frame}

% Depuración
\begin{frame}[fragile]{Depuración}
    \begin{itemize}
        \item Curva de aprendizaje pronunciada.
        
        \item La meta es tener un programa libre de errores.
        
        \item {\blue\bfseries Herramientas}
        \begin{itemize}
            \item {\blue\bfseries Cerebro}: ayuda a cazar {\blue sistemáticamente} el error.
            \item Módulo \lstinline|pdb|.
            \item Visual Studio Code, spyder.
            \item Instrucciones \lstinline|print| (¡no abusar!)
        \end{itemize}
    \end{itemize}
\end{frame}

% Depuración con print
\begin{frame}[fragile]{Depuración con \lstinline|print|}
    \begin{itemize}
        \item Buena forma de probar hipótesis.
        
        \item {\blue\bfseries ¿Cuándo?}
        \begin{itemize}
            \item Entrar a funciones.
            \item Parámetros
            \item Resultados intermedios, finales. 
        \end{itemize}
        
        \item {\blue\bfseries Método de bisección}
        \begin{itemize}
            \item \lstinline|print| al centro del código.
            \item Los errores dan una idea de {\blue dónde} puede estar el error.
        \end{itemize}
    \end{itemize}
\end{frame}

\begin{frame}{}
    \begin{figure}
        \centering
        \includegraphics[width=0.6\textwidth]{figs/depuracion/debugging-patrick-meme.jpg}
    \end{figure}
\end{frame}

% Pasos de depuración
\begin{frame}[fragile]{Pasos de depuración}
    \begin{itemize}
        \item {\orange\bfseries Estudiar} el código fuente:
        \begin{itemize}
            \item ¿Cómo obtuve el resultado inesperado?
            \item ¿Es parte de alguna familia de erores?
        \end{itemize}
        
        \item {\orange\bfseries Método científico}:
        \begin{itemize}
            \item Estudiar los datos disponibles.
            \item Formar una hipótesis.
            \item Experimento.
            \item Escoger una entrada simple para probar.
        \end{itemize}
    \end{itemize}
\end{frame}

% Familias de errores
\includeslides{15}{slides-60001/MIT6_0001F16_Lec7_mod.pdf}{Familias de errores}

% Errores de lógica
\includeslides{16}{slides-60001/MIT6_0001F16_Lec7_mod.pdf}{Errores de lógica : difíciles}

% ¿Qué no hacer?, ¿qué sí?
\includeslides{17}{slides-60001/MIT6_0001F16_Lec7_mod.pdf}{¿Qué no hacer?, ¿qué sí?}



\section{Manejo de excepciones}

% Excepciones
\includeslides{18-19}{slides-60001/MIT6_0001F16_Lec7_mod.pdf}{Excepciones}

% Manejo de excepciones
\begin{frame}[fragile]{Manejo de excepciones}
    \begin{itemize}
        \item Python provee bloques de {\blue\bfseries manejo} de excepciones (exceptions).
\begin{lstlisting}
try:
    a = int(input("Tell me one number:"))
    b = int(input("Tell me another number:"))
    print(a/b)
except:
    print("Bug in user input.")
\end{lstlisting}

        \item Las {\orange\bfseries excepciones} que surjan (\textit{raised by}) en el cuerpo de una sentencia \lstinline|try| son manejadas por el cuerpo de \lstinline|except| mientras {\blue la ejecución continúa}.
    \end{itemize}
\end{frame}

% Excepciones específicas
\begin{frame}[fragile]{Excepciones específicas}
    \begin{itemize}
        \item Es posible especificar cláusulas de \lstinline|except| para una excepción en particular. 
\begin{lstlisting}
try:
    a = int(input("Tell me one number: "))
    b = int(input("Tell me another number: "))
    print("a/b = ", a/b)
    print("a+b = ", a+b)
except ValueError:
    print("Could not convert to a number.")
except ZeroDivisionError:
    print("Can't divide by zero")
except:
    print("Something went very wrong.")
\end{lstlisting}

    \end{itemize}
\end{frame}

% Más cláusulas
\begin{frame}[fragile]{Otras excepciones}
    \begin{itemize}
        \item \lstinline|else|:
        \begin{itemize}
            \item Se ejecuta si la ejecución de \lstinline|try| {\orange se completa sin excepciones}.
        \end{itemize}
        
        \item \lstinline|finally|:
        \begin{itemize}
            \item Se ejecuta {\orange\bfseries siempre} después de \lstinline|try|, \lstinline|else| y \lstinline|except| o si ejecuta \lstinline|break|, \lstinline|continue| o \lstinline|return|.
            
            \item {\orange\bfseries Útil} para limpiar el código con instrucciones que de todas formas deben ejecutarse sin importar el resultado, como cerrar un archivo, conexión HTTP, etc.
        \end{itemize}
    \end{itemize}
\end{frame}

% ¿Qué hacer con las excepciones?
\begin{frame}[fragile]{¿Qué hacer con las excepciones?}
    \begin{itemize}
        \item {\blue\bfseries Manejarlas silenciosamente}
        \begin{itemize}
            \item Sustituir valores por defecto, solamente continuar con el código.
            \item ¡Mala idea! el usuario no tiene advertencias.
        \end{itemize}
        
        \item Devolver un valor para {\blue\bfseries representar} el error. 
        \begin{itemize}
            \item ¿Qué valor escoger?
            \item Complica el resto del código.
        \end{itemize}
        
        \item Detener la ejecución, {\orange\bfseries mostrar una condición de error}:
        \begin{itemize}
            \item En Python: \lstinline|raise|
        \end{itemize}
\begin{lstlisting}
raise Exception("string descriptivo")
\end{lstlisting}
    \end{itemize}
\end{frame}

% Exception example
\begin{frame}[fragile]{Excepciones: ejemplo}
\begin{lstlisting}
def get_ratios(L1, L2):
    """ Assumes: L1, L2 lists of equal length of numbers
        Returns: a list containing L1[i]/L2[i] """
    ratios = []
    for index in range(len(L1)):
        try:
            ratios.append(L1[index]/L2[index])
        except ZeroDivisionError:
            ratios.append(float('nan'))
            #nan = not a number
        except:
            raise ValueError('called with bad arg')
    return ratios
\end{lstlisting}
\end{frame}

%% Ejemplos adicionales de try-except
\includeslides{26-30}{slides-60001/MIT6_0001F16_Lec7_mod.pdf}{Excepciones: otro ejemplo}


\subsection{Assertions}

\begin{frame}[fragile]{Uso de \lstinline|assert|}
    \begin{itemize}
        \item Cuando quieres estar {\blue\bfseries seguro} de que los supuestos del cómputo son como se esperaba.
        
        \item Utilizamos una cláusula \lstinline|assert| para generar una excepción \lstinline|AssertionError| si no se verifican los supuestos.
        
        \item Este es un ejemplo de buena {\blue\bfseries programación defensiva}.
    \end{itemize}
\end{frame}

% Ejemplo de assert
\begin{frame}[fragile]{Ejemplo}
\begin{lstlisting}
def avg(grades):
    """ grades : grades list """
    assert len(grades) != 0, 'no grades data'
    return sum(grades)/len(grades)
\end{lstlisting}

    \begin{itemize}
        \item Esto lleva a una excepción de tipo \lstinline|AssertionError| si se da una lista vacía de notas.
        
        \item De lo contrario, ¡OK!
    \end{itemize}
\end{frame}

% Assertions como programación defensiva
\begin{frame}[fragile]{\textit{Assertions} como programación defensiva}
    \begin{itemize}
        \item Las afirmaciones no permiten que el programador controle respuesta a condiciones inesperadas.
        
        \item Garantizan que la {\blue ejecución se detenga} siempre que condición no se cumple.
        
        \item Normalmente se usan para {\blue verificar las entradas} a las funciones, pero pueden ser usadas en cualquier lugar.
        
        \item Pueden usarse para {\blue verificar las salidas} de una función para evitar propagar malos valores.
        
        \item Pueden facilitar la depuración de un error.
    \end{itemize}
\end{frame}

\begin{frame}[fragile]{¿Dónde usarlas?}
    \begin{itemize}
        \item El objetivo es encontrar errores tan pronto como se introduzcan.
        
        \item {\blue\bfseries Suplemento} a la etapa de pruebas.
        
        \item Levantar {\blue\bfseries excepciones} si el usuario provee datos incorrectos.
        
        \item También permiten: 
        \begin{itemize}
            \item Verificar {\blue\bfseries tipos} de los argumentos o valores.
            \item Revisar {\blue\bfseries anomalías} en las estructuras de datos.
            \item Revisar {\blue\bfseries restricciones} en los valores de retorno.
            \item Revisar {\blue\bfseries violaciones} en las restricciones de un procedimiento.
        \end{itemize}
    \end{itemize}
\end{frame}


%-----------------------------------------------------------------
% Termina el contenido de la presentación
%-----------------------------------------------------------------

\begin{comment}

% Slide with two columns, half and half
\begin{frame}[fragile]{}
    \begin{columns}
        \begin{column}{0.5\textwidth}
        \end{column}
        \begin{column}{0.5\textwidth}
        \end{column}
    \end{columns}
\end{frame}

% Slide with bullets
\begin{frame}[fragile]{}
    \begin{itemize}
        \item 
    \end{itemize}
\end{frame}

% Slide with bullets and subbullets
\begin{frame}[fragile]{}
    \begin{itemize}
        \item {\blue\bfseries A}
        \begin{itemize}
            \item a
        \end{itemize}
    \end{itemize}
\end{frame}

% Slide with bullets and code
\begin{frame}[fragile]{}
    \begin{itemize}
        \item 
    \end{itemize}
\begin{lstlisting}
Code goes here
\end{lstlisting}
\end{frame}

% Slides PDF
\includeslides{54-57}{slides-60001/MIT6_0001F16_Lec6_mod.pdf}{Eficiencia: Fibonacci con diccionarios}

% Pytutor
\href{URL}{Ver en Python Tutor}


\end{comment}

\end{document}